% ---
% id: foundations/req:declared-tolerance
% kind: req
% title: "One declared tolerance per instance"
% summary: >-
%   One tolerance in [0, 1) per instance, with no default, caps the defect,
%   measured in binary64, that an instance may repair in what it declares. A
%   family snapped on an exact reference has no defect, supplier output is
%   decided by its consumer, a prepared probability vector by its own rule.
% prominence: standard
% status: accepted
% uses: [def:tolerance, def:snap, def:declared-data, def:exact-reference]
% mentions: [req:who-snaps, def:exact-island, req:boundary-admission, req:snap-reference]
% assumes: []
% refs: []
% provenance: newly-derived
% notes: "A family whose snap runs on the data as they stand, an evaluation of a model with an enclosed reference among them, falls under the first paragraph, measured against tau like data without a reference; the second paragraph covers a snap that runs on the centres of an exact reference."
% lean:
% ---

The domain stops below one because at a tolerance of one or more a
completeness test would admit the all-zero family, whose completeness residue
against the identity is exactly one on every block. How an instance turns
ordinary numerical inputs into \tref[declared value]{declared values}, the
representation policy by which a binary64 input denotes a rounding cell
included, is documented with the run tier
(\repodoc{runs/input-preparation.md}).

\begin{requirement}\label{req:declared-tolerance}
Each instance declares exactly one \tref{declared tolerance} $\tnum$
(Definition~\ref{def:tolerance}) in its run configuration, and $\tnum$ has no
default: a run whose configuration does not declare it, or declares a value
that is not a binary64 number in $[0,1)$, is refused before anything is
constructed. $\tnum$ is the ceiling on the defect an instance may repair in
what it declares. It binds the data an instance is about to declare other than
numbers of a supplier's output that it passes on unchanged
(Definition~\ref{def:declared-data}\,(ii)); those it neither measures nor
repairs, and their consumer decides their conditions exactly on entry
(Requirement~\ref{req:boundary-admission}). Of the data it binds that have no
\tref{exact reference} on whose centres the snap runs (for those, see the
next paragraph), the instance
measures in binary64 the defects (Definition~\ref{def:tolerance}) and, for each
block read as positive semidefinite, the least eigenvalue of its Hermitian
part, each on the data as they stand before any step of the snap, and records
those measured numbers. Data with a measured defect above $\tnum$ or a
measured least eigenvalue below $-\tnum$ are refused as a modelling error; the
others may be snapped (Definition~\ref{def:snap}), the snap being recorded
with the measured defects it corrected (Requirement~\ref{req:who-snaps}). Near
$\tnum$ this decision on measured numbers may differ from the predicate
\emph{\tref{within the declared tolerance}} of Definition~\ref{def:tolerance},
whose sums are exact, and no bound on that difference is kept. This decision
is not on the \tref{exact island}; the exact decisions are those by which the
snap establishes the conditions of its output
(Definition~\ref{def:exact-island}\,(i)).

An operator family that an instance snaps on the centres of an
\tref{exact reference} (Definition~\ref{def:exact-reference}), as
Requirement~\ref{req:snap-reference}\,(ii) states, stands before the snap as
that reference, which is Hermitian, positive semidefinite and complete by the
exact decision on its entries, and never by a binary64 evaluation. Its defects are
therefore zero and no block of it falls below the \tref{positivity floor}:
$\tnum$ decides nothing for it, and the instance records the defects as zero
on that ground rather than on a binary64 measurement. $\tnum$ keeps its role
for the data it binds on which the snap runs as they stand, an evaluation of
a model whose reference has a radius that is not zero included, since the
snap does not run on such a reference.

A probability vector that an instance prepares from ordinary numerical inputs
before declaring it is decided instead by the rule of the input preparation.
It is refused unless its normalisation defect, computed exactly on the inputs
as given, is at most $\tnum$, or the sum of the rounding cells of its inputs
contains $1$: a binary64 input contributes the rounding cell it denotes under
the preparation's documented representation policy, and an input given
exactly, as a decimal or a fraction, contributes the one point that is its
value. A value below $0$ or above $1$ is refused. The preparation records the
defect, the cells and which of the two grounds admitted the vector. Its snap
then runs at a \emph{guard}, the larger of $\tnum$ and the completeness
residue the snap measures on the binary64 conversion of the inputs, which may
exceed $\tnum$; the guard is recorded, and it admits no vector the rule has
refused.
\end{requirement}
